\documentclass[11pt]{article}
\usepackage{amsmath,amsthm,amssymb}
\usepackage[margin=1in]{geometry}
\usepackage{enumitem}
\usepackage{booktabs}

\newtheorem{theorem}{Theorem}
\newtheorem{lemma}[theorem]{Lemma}
\newtheorem{proposition}[theorem]{Proposition}
\newtheorem{corollary}[theorem]{Corollary}
\newtheorem{conjecture}[theorem]{Conjecture}
\theoremstyle{definition}
\newtheorem{definition}[theorem]{Definition}
\newtheorem{remark}[theorem]{Remark}
\newtheorem{example}[theorem]{Example}

\title{The $j$-formula for the family $\lambda=(2,2,1^{n-4})$:\\
  the first non-hook structural proof}
\author{Clio}
\date{13 May 2026}

\begin{document}
\maketitle

\begin{abstract}
We extend the May-11 and May-12 structural proofs of the $j$-formula
from hooks $\lambda=(2,1^{n-2})$ and $\lambda=(3,1^{n-3})$ to the
first non-hook family in the rank-zero regime: $\lambda=(2,2,1^{n-4})$
at every $n \ge 6$.  The proof tracks the image of the staircase
product $R'_{n-1}R'_n$ on $V_\lambda$ step-by-step in the seminormal
basis.  After $R'_n$ the image is exactly the $+q$-eigenspace of
$T_{n-1}$, of dimension $n-3$, with explicit basis of $T_{n-1}$
$2$-block $+q$-eigenvectors.  The first factor $(T_1{+}1)$ of
$R'_{n-1}$ collapses this to a $1$-dimensional subspace; subsequent
factors shift this subspace through a $4$-element ``window'', and the
final factor $(T_{n-2}{+}1)$ produces a 1-dimensional image with
support on exactly $4$ standard tableaux, all in $E_1^-$.  As a
corollary, $\Pi^{S_n}|_{V_{(2,2,1^{n-4})}}=0$ unconditionally for
every $n\ge 6$.

This is the first non-hook case where the rank-zero theorem is
established structurally.  Combined with the hook results of May 11
and May 12, the rank-zero theorem now holds unconditionally for three
infinite families in the rank-zero regime.
\end{abstract}

\section{Setup and notation}

We work with the Hecke algebra $H_q(S_n)$ over $\mathbb{Q}(q)$,
with generators $T_1,\ldots,T_{n-1}$ satisfying
$T_i^2 = (q-1)T_i + q$ and the standard braid relations.  $V_\lambda$
denotes the seminormal cell module of $H_q(S_n)$ for the partition
$\lambda \vdash n$.  Its basis $\{v_T\}$ is indexed by standard Young
tableaux (SYTs) of shape $\lambda$.

The staircase Bott--Samelson element is
\[
   \Pi^{S_n} = R'_2 R'_3 \cdots R'_n,\qquad
   R'_k = (T_{k-1}{+}1)(T_{k-2}{+}1)\cdots(T_1{+}1).
\]
Write $B_j := R'_j R'_{j+1}\cdots R'_n$.  $E_i^\pm$ denotes the
$\pm$-eigenspace of $T_i$ (eigenvalues $q$ resp.\ $-1$).

\subsection{The seminormal basis of $\lambda=(2,2,1^{n-4})$}

Throughout, fix $\lambda = (2,2,1^{n-4})$ for some $n \ge 6$.  The
shape has length $\ell(\lambda) = n-2$, total $|\lambda| = n$.

\begin{lemma}[SYTs of $\lambda=(2,2,1^{n-4})$]\label{lem:syts-22}
SYTs of shape $\lambda = (2,2,1^{n-4})$ are in bijection with pairs
$(a,c)$ satisfying $2 \le a < c \le n$ and $(a,c) \ne (2,3)$.
\end{lemma}

\begin{proof}
For an SYT $T$ of shape $\lambda$, write $a$ for the entry at $(1,2)$
and $c$ for the entry at $(2,2)$.  Necessarily $a > 1$, $c > a$, and
the entry $b$ at $(2,1)$ is the smallest element of
$\{2,\ldots,n\}\setminus\{a,c\}$.  Validity of the SYT requires
$b < c$; this fails iff $a = 2,\, c = 3$ (in which case $b = 4 > 3$),
and otherwise holds.  All remaining entries fill the rest of column~1
in increasing order.
\end{proof}

We write $v_{(a,c)} := v_T$ for $T$ the SYT with second-column
entries $(a,c)$.  We refer to ``column~1'' meaning the leftmost column
(positions $(1,1),(2,1),\ldots,(n-2,1)$) and ``column~2'' meaning
positions $(1,2),(2,2)$.

\begin{lemma}[$T_i$-action on the seminormal basis of $\lambda$]\label{lem:Ti-action}
For $i \in \{1,\ldots,n-1\}$ and SYT $T = T_{(a,c)}$:
\begin{itemize}[leftmargin=2em,topsep=0.3em,itemsep=0.2em]
\item If $i,i{+}1$ lie in the same row of $T$: $T_iv_{(a,c)} = qv_{(a,c)}$.
\item If $i,i{+}1$ lie in the same column of $T$: $T_iv_{(a,c)} = -v_{(a,c)}$.
\item Otherwise (different rows and columns), $T_i$ acts as a
$2\times 2$ Hoefsmit block on $\mathrm{span}(v_T, v_{s_iT})$ where
$s_i T$ is the SYT obtained by swapping $i$ and $i{+}1$.  At axial
distance $d := c(i{+}1) - c(i)$ (where $c(r,k) = k - r$), the
$+q$-eigenvector is, up to scalar,
\[
   u^+_{T,s_iT} = \alpha_d v_T + \beta_d v_{s_iT},
   \qquad
   (T_i{+}1)v_T = \kappa_d\, u^+_{T,s_iT},
\]
with $\alpha_d, \beta_d, \kappa_d \in \mathbb{Q}(q)$ generically nonzero.
\end{itemize}
\end{lemma}

\begin{proof}
Standard: see e.g.\ Hoefsmit's formula or Murphy's seminormal basis.
\end{proof}

\begin{lemma}[$E_1^\pm$ for $\lambda=(2,2,1^{n-4})$]\label{lem:E1-22}
$T_1$ acts diagonally on $V_\lambda$ (no $2$-blocks).  Specifically:
\[
   E_1^+ = \mathrm{span}\{v_{(2,c)} : 4 \le c \le n\},
   \qquad
   E_1^- = \mathrm{span}\{v_{(a,c)} : 3 \le a < c \le n\}.
\]
\end{lemma}

\begin{proof}
$T_1v_{(a,c)}$ depends on the positions of $1$ and $2$.  $1$ is always
at $(1,1)$.  $2$ is at $(1,2)$ iff $a = 2$ (same row as $1$,
$T_1v=qv$), and otherwise at $(2,1)$ (same column as $1$, $T_1v=-v$).
The two cases partition $V_\lambda$.
\end{proof}

\section{The main theorem}

\begin{theorem}[$j$-formula for $\lambda=(2,2,1^{n-4})$]\label{thm:main}
For every $n \ge 6$, $\lambda = (2,2,1^{n-4}) \vdash n$ has length
$\ell = n-2$ and the $j$-formula holds with $j = \ell + 1 = n-1$:
\[
   B_{n-1}(V_\lambda) := R'_{n-1}\,R'_n\,V_\lambda \;\subseteq\; E_1^-.
\]
Moreover, the image is one-dimensional, with support exactly on the
$4$ SYTs indexed by
\[
   \bigl\{(n{-}3,\,n{-}2),\ (n{-}3,\,n{-}1),\ (n{-}2,\,n),\ (n{-}1,\,n)\bigr\}.
\]
\end{theorem}

\begin{corollary}[$\Pi^{S_n}|_{V_{(2,2,1^{n-4})}} = 0$]\label{cor:rank-zero}
$\Pi^{S_n}|_{V_{(2,2,1^{n-4})}} = 0$ unconditionally for every
$n \ge 6$.
\end{corollary}

\begin{proof}
By Theorem~\ref{thm:main}, $B_{n-1}V_\lambda \subseteq E_1^-$.  The
factorization $\Pi^{S_n} = R'_2 \cdots R'_{n-2} \cdot B_{n-1}$ has
$R'_{n-2}$ ending in $(T_1{+}1)$, which annihilates $E_1^-$.  Hence
$\Pi^{S_n} V_\lambda = 0$.
\end{proof}

The proof of Theorem~\ref{thm:main} occupies
Sections~\ref{sec:phaseA} (analysis of $R'_n$) and \ref{sec:phaseB}
(analysis of $R'_{n-1}$ on the result of $R'_n$).

\section{Phase A: $R'_n$ surjects onto $E_{n-1}^+$}\label{sec:phaseA}

For $1 \le j \le n-1$, write
$W_j^{(n)} := (T_j{+}1)(T_{j-1}{+}1)\cdots(T_1{+}1)\,V_\lambda$.
Note $W_{n-1}^{(n)} = R'_n V_\lambda$.

\begin{proposition}[Phase A endpoint]\label{prop:phaseA}
$W_j^{(n)} = E_j^+$ for every $1 \le j \le n-1$.  In particular,
$\dim W_j^{(n)} = n-3$.
\end{proposition}

The proof proceeds via an explicit description of a basis of
$W_j^{(n)}$.

\begin{definition}[Basis $\mathcal{B}_j$]\label{def:Bj}
For $1 \le j \le n-1$, define $\mathcal{B}_j$ to be the set of vectors
in $V_\lambda$ given by:
\begin{itemize}[leftmargin=2em,topsep=0.3em,itemsep=0.2em]
\item For each $a \in \{2,\ldots,j-1\}$ such that $(a,j)$ is a valid
SYT label (i.e., $(a,j) \ne (2,3)$), the $T_j$ $+q$-eigenvector
$w_j(a)$ in the $2$-block $\{v_{(a,j)}, v_{(a,j+1)}\}$.  When the
$2$-block degenerates because $(a,j)=(2,3)$ is invalid (i.e., $a=2,
j=3$), we set $w_3(2) := v_{(2,4)}$ (the unique row-$2$ same-row
$+q$-eigenvector of $T_3$ in $V_\lambda$).
\item For each $c \in \{j+2,\ldots,n\}$, the $T_j$ $+q$-eigenvector
$w_j(c)$ in the $2$-block $\{v_{(j,c)}, v_{(j+1,c)}\}$.
\end{itemize}
\end{definition}

\begin{remark}[Block validity]\label{rem:blocks}
We check that the $2$-blocks above genuinely exist.  The block
$\{v_{(a,j)}, v_{(a,j+1)}\}$ requires $T_j$ to couple them, which
happens when $j$ and $j{+}1$ are in different rows and columns of
$T_{(a,j)}$: $j = c$ at $(2,2)$, while $j+1$ is in column~$1$
(verifiable when $a < j$ and $(a,j) \ne (2,3)$).  The block
$\{v_{(j,c)}, v_{(j+1,c)}\}$ requires $j = a$ at $(1,2)$ and $j+1$ in
column~$1$, valid when $j+1 < c$.  All cases are checked in the
inductive proof below.
\end{remark}

\begin{lemma}[Cardinality of $\mathcal{B}_j$]\label{lem:Bj-size}
$|\mathcal{B}_j| = n - 3$ for every $1 \le j \le n - 1$.
\end{lemma}

\begin{proof}
Counting: the $a$-indexed contribute $j - 2$ vectors (one per $a \in
\{2,\ldots,j-1\}$, with the degenerate case $j=3, a=2$ still
counted as one vector), and the $c$-indexed contribute $n - j - 1$
vectors (one per $c \in \{j+2,\ldots,n\}$).  Total: $(j-2) + (n-j-1) =
n-3$.  For $j = 1, 2$, the $a$-indexed range is empty, and the
$c$-indexed range $\{j+2,\ldots,n\}$ has $n - j - 1$ elements
intersected with the constraint $(j,c)$ or $(j+1,c)$ are valid SYT
labels --- we verify case-by-case below.

For $j = 1$: $\mathcal{B}_1 = \{w_1(c) : c \in \{3,\ldots,n\}\}$ where
$w_1(c)$ is supposed to be the $+q$-eigenvector in the block
$\{v_{(1,c)}, v_{(2,c)}\}$.  Since $(1,c)$ is not a valid SYT label,
there is no $2$-block; instead $T_1$ acts diagonally and $w_1(c) :=
v_{(2,c)}$ for $c \ge 4$ (with $c = 3$ excluded as $(2,3)$ is not a
valid SYT label).  This gives $|\mathcal{B}_1| = n - 3$.

For $j = 2$: $\mathcal{B}_2 = \{w_2(c) : c \in \{4,\ldots,n\}\}$.
$|\mathcal{B}_2| = n - 3$.
\end{proof}

\begin{lemma}[Inductive transition $\mathcal{B}_{j-1} \to \mathcal{B}_j$]\label{lem:transition}
For $2 \le j \le n-1$, the map $(T_j{+}1)$ takes $\mathcal{B}_{j-1}$
bijectively (up to nonzero scalars) to $\mathcal{B}_j$, via the
transformation:
\begin{itemize}[leftmargin=2em,topsep=0.3em,itemsep=0.2em]
\item $w_{j-1}(a) \;\mapsto\; (\text{nonzero scalar}) \cdot w_j(a)$ for
$a \in \{2,\ldots,j-2\}$ (those $a$-indexed vectors of
$\mathcal{B}_{j-1}$ for which $a \le j-2$).
\item $w_{j-1}(c=j+1) \;\mapsto\; (\text{nonzero scalar}) \cdot w_j(a=j-1)$
(this maps the $c$-indexed vector with $c = j$ in $\mathcal{B}_{j-1}$,
i.e., $c = (j-1)+2$, to the new $a$-indexed vector at
$a = j-1$ in $\mathcal{B}_j$).
\item $w_{j-1}(c) \;\mapsto\; (\text{nonzero scalar}) \cdot w_j(c)$
for $c \in \{j+2,\ldots,n\}$.
\end{itemize}
In particular, $\mathcal{B}_j$ is a basis of $W_j^{(n)}$, so
$\dim W_j^{(n)} = n-3 = \dim E_j^+$, hence $W_j^{(n)} = E_j^+$.
\end{lemma}

\begin{proof}
We compute $(T_j{+}1)$ applied to each basis vector of
$\mathcal{B}_{j-1}$ in the seminormal basis.

\paragraph{(I) $w_{j-1}(a)$ for $a \in \{2,\ldots,j-2\}$.}
By Definition~\ref{def:Bj} (and assuming $a < j-1$, $(a, j-1) \ne (2,3)$),
$w_{j-1}(a) = \alpha v_{(a,j-1)} + \beta v_{(a,j)}$.

\emph{Compute $(T_j{+}1)v_{(a,j-1)}$.} In $T_{(a,j-1)}$ with
$a \le j-2$: $a$ at $(1,2)$, $c = j-1$ at $(2,2)$.  $j$ and $j+1$ are
both in column~$1$ (since neither equals $a$ nor $c$).
Counting: column~$1$ entries are $\{1\} \cup \{2,\ldots,n\}\setminus\{a,j-1\}$;
elements $\le j-1$ in this set total $j-2$, so position of $j$ is row
$j-1$ if $j$ is the $(j-1)$th smallest.  Direct verification: column~$1$
entries less than $j$ are $\{1, 2, \ldots, j-1\} \setminus \{a, j-1\}$,
totaling $j - 3$ elements, plus $\{1\}$ gives $j-2$.
Therefore $j$ is at row $j-1$, position $(j-1,1)$.
Likewise $j+1$ at $(j,1)$.  Same column adjacent: $T_jv = -v$, so
$(T_j{+}1)v_{(a,j-1)} = 0$.

\emph{Compute $(T_j{+}1)v_{(a,j)}$.} In $T_{(a,j)}$: $a$ at $(1,2)$,
$c = j$ at $(2,2)$.  $j+1$ is in column~$1$ (not equal to $a$ or
$c$).  Counting: column~$1$ entries less than $j+1$ are
$\{1,\ldots,j\}\setminus\{a,j\}$, totaling $j-2$; so $j+1$ is at row
$j-1$, position $(j-1,1)$, content $-(j-2)$.
The cell of $j$ is $(2,2)$, content $0$.  Different positions.
$T_j$-block partner: swap $j,j+1$ moves $j+1$ to $(2,2)$, $j$ to
$(j-1,1)$; new SYT label $(a, j+1)$.
Axial distance $d = -(j-2) - 0 = -(j-2)$, so
$(T_j{+}1)v_{(a,j)} = \kappa\,u^+_{T_{(a,j)}, T_{(a,j+1)}}
= \kappa(\alpha' v_{(a,j)} + \beta' v_{(a,j+1)})$
for nonzero $\alpha',\beta',\kappa$ depending on $j,n,q$.

Combining: $(T_j{+}1)w_{j-1}(a) = \beta\kappa(\alpha' v_{(a,j)} +
\beta' v_{(a,j+1)})$, which is a nonzero scalar multiple of $w_j(a)$.
\checkmark

\paragraph{(II) $w_{j-1}(c=j+1)$.}
By Definition~\ref{def:Bj}, $w_{j-1}(c=j+1) = \alpha v_{(j-1,j+1)} +
\beta v_{(j,j+1)}$ ($T_{j-1}$-block).

\emph{Compute $(T_j{+}1)v_{(j-1,j+1)}$.} In $T_{(j-1,j+1)}$: $a = j-1$
at $(1,2)$, $c = j+1$ at $(2,2)$.  $j$ in column~$1$ (not equal to
$a$ or $c$).  Position: $j$ has $j-2$ smaller elements in column~$1$
($\{1, 2, \ldots, j-1\} \setminus \{j-1\} = \{1,\ldots,j-2\}$), so
$j$ at $(j-1,1)$, content $-(j-2)$.  Different from $(2,2)$ where
$j+1$ sits.  $T_j$-block partner: swap $j,j+1$ moves $j+1$ to
$(j-1,1)$, $j$ to $(2,2)$; new SYT $T_{(j-1,j)}$.
Axial distance $0 - (-(j-2)) = j-2$, so
$(T_j{+}1)v_{(j-1,j+1)} \propto u^+ = \alpha'v_{(j-1,j)} +
\beta'v_{(j-1,j+1)}$ (a $T_j$ $2$-block $+q$-eigenvector with support
on $\{(j-1,j),(j-1,j+1)\}$).

\emph{Compute $(T_j{+}1)v_{(j,j+1)}$.} $a = j, c = j+1$: same column~$2$.
$T_jv = -v$, so $(T_j{+}1)v = 0$.

Combining: $(T_j{+}1)w_{j-1}(c=j+1) = \alpha\kappa(\alpha' v_{(j-1,j)}
+ \beta' v_{(j-1,j+1)})$, a nonzero scalar multiple of $w_j(a=j-1)$.
\checkmark

\paragraph{(III) $w_{j-1}(c)$ for $c \in \{j+2,\ldots,n\}$.}
$w_{j-1}(c) = \alpha v_{(j-1,c)} + \beta v_{(j,c)}$.

\emph{Compute $(T_j{+}1)v_{(j-1,c)}$.} $a = j-1, c \ge j+2$.
$j$ and $j+1$ both in column~$1$ (since neither equals $a$ nor $c$).
Adjacent rows.  $T_jv = -v$, so $(T_j{+}1)v = 0$.

\emph{Compute $(T_j{+}1)v_{(j,c)}$.} $a = j$ at $(1,2)$, content $1$.
$j+1$ in column~$1$ (not equal to $a$, and $c \ge j+2 \ne j+1$).
Counting: column~$1$ entries less than $j+1$ are
$\{1, 2, \ldots, j\} \setminus \{j\} = \{1,\ldots,j-1\}$, totaling
$j-1$.  So $j+1$ is at $(j,1)$, content $1-j$.  Different from
$(1,2)$.  $T_j$-block partner: swap moves $j+1$ to $(1,2)$, $j$ to
$(j,1)$; new SYT $T_{(j+1,c)}$.  Axial distance $1-j-1=-j$.
$(T_j{+}1)v_{(j,c)} \propto u^+ = \alpha'v_{(j,c)} + \beta'v_{(j+1,c)}$,
a $T_j$ $2$-block $+q$-eigenvector with support $\{(j,c),(j+1,c)\}$.

Combining: $(T_j{+}1)w_{j-1}(c) = \beta\kappa(\alpha'v_{(j,c)}
+ \beta'v_{(j+1,c)})$, a nonzero scalar multiple of $w_j(c)$.
\checkmark

\paragraph{Conclusion.}
$(T_j{+}1)$ takes $\mathcal{B}_{j-1}$ to a set of $n-3$ vectors, each
a nonzero scalar multiple of a distinct element of $\mathcal{B}_j$.
The $n-3$ vectors of $\mathcal{B}_j$ have pairwise distinct supports
(by inspection: the $a$-indexed have support
$\{(a,j),(a,j+1)\}$ with $a \le j-1$, the $c$-indexed have support
$\{(j,c),(j+1,c)\}$ with $c \ge j+2$, and the supports are pairwise
disjoint).  Hence $\mathcal{B}_j$ is linearly independent, and spans
$W_j^{(n)}$.

The image $W_j^{(n)} \subseteq E_j^+$ since $(T_j{+}1) = (1+q)\pi_j^+$
where $\pi_j^+$ is the projection onto $E_j^+$ along $E_j^-$.
Comparing dimensions: $|\mathcal{B}_j| = n-3$ and we'll show
$\dim E_j^+ \le n-3$ by counting (Lemma~\ref{lem:Ej-dim} below), giving
equality.
\end{proof}

\begin{lemma}[Dimension of $E_j^+$]\label{lem:Ej-dim}
For each $1 \le j \le n-1$, $\dim E_j^+ = n - 3$.
\end{lemma}

\begin{proof}
Since $T_j$ and $T_{j'}$ are conjugate in $H_q(S_n)$ (any two simple
reflections are conjugate via a sequence of braid moves), their traces
on the irreducible $V_\lambda$ are equal.  Hence
$\mathrm{tr}(T_j|_{V_\lambda})$ is a constant in $j$.  As $T_j$ is
diagonalizable on $V_\lambda$ with eigenvalues $\{q, -1\}$,
\[
   \mathrm{tr}(T_j|_{V_\lambda})
   = q\cdot\dim E_j^+ - \dim E_j^-,
   \qquad
   \dim E_j^+ + \dim E_j^- = \dim V_\lambda = \tfrac{n(n-3)}{2}.
\]
Solving, $\dim E_j^+ = (\mathrm{tr}(T_j) + \dim V_\lambda)/(q + 1)$,
which is independent of $j$.

Compute at $j = 1$ using Lemma~\ref{lem:E1-22}: $\dim E_1^+ = n - 3$
(the SYTs $v_{(2,c)}$ for $4 \le c \le n$).  Hence $\dim E_j^+ = n-3$
for all $j$.
\end{proof}

\begin{proof}[Proof of Proposition~\ref{prop:phaseA}]
Combining Lemmas~\ref{lem:Bj-size}, \ref{lem:transition},
\ref{lem:Ej-dim}: $W_j^{(n)} \subseteq E_j^+$, and both have dimension
$n-3$, so $W_j^{(n)} = E_j^+$.
\end{proof}

\begin{corollary}[$W_n$ described]\label{cor:Wn}
$W_n := R'_n V_\lambda = E_{n-1}^+$ has basis
$\mathcal{B}_{n-1} = \{w_{n-1}(a) : a \in \{2,\ldots,n-2\}\}$
where $w_{n-1}(a) = \alpha v_{(a,n-1)} + \beta v_{(a,n)}$ is the
$T_{n-1}$ $+q$-eigenvector in the $2$-block
$\{v_{(a,n-1)}, v_{(a,n)}\}$.
\end{corollary}

\begin{proof}
At $j = n-1$: $a$-indexed range is $\{2,\ldots,n-2\}$ ($n-3$ elements),
$c$-indexed range $\{n+1,\ldots,n\}$ is empty.  Each $w_{n-1}(a)$ is
a $T_{n-1}$ $2$-block $+q$-eigenvector by construction.
\end{proof}

\section{Phase B: collapse to $E_1^-$ via $R'_{n-1}$}\label{sec:phaseB}

We now apply $R'_{n-1} = (T_{n-2}{+}1)\cdots(T_1{+}1)$ to $W_n$.
Define $\widetilde{W}_j := (T_j{+}1)\cdots(T_1{+}1)\,W_n$ for
$1 \le j \le n-2$.

\begin{proposition}[Step $1$: collapse to dim $1$]\label{prop:step1}
$\widetilde{W}_1 = \mathrm{span}(w_{n-1}(2))$, $1$-dimensional.
\end{proposition}

\begin{proof}
By Corollary~\ref{cor:Wn}, $W_n$ has basis $\{w_{n-1}(a) : 2 \le a \le n-2\}$
with $w_{n-1}(a) = \alpha v_{(a,n-1)} + \beta v_{(a,n)}$.

\emph{Case $a = 2$.} Both $v_{(2,n-1)}$ and $v_{(2,n)}$ have
$2 \in \{a\}$ in row~$1$, so by Lemma~\ref{lem:E1-22} they lie in
$E_1^+$.  Hence $T_1 w_{n-1}(2) = q\,w_{n-1}(2)$, so
$(T_1{+}1) w_{n-1}(2) = (1+q)\,w_{n-1}(2)$.  Survives.

\emph{Case $a \ge 3$.} Both $v_{(a,n-1)}$ and $v_{(a,n)}$ have $a \ge 3$,
so $2$ is in column~$1$, hence in $E_1^-$ (Lemma~\ref{lem:E1-22}).
$(T_1{+}1) w_{n-1}(a) = 0$.  Killed.

Therefore $\widetilde{W}_1 = \mathrm{span}(w_{n-1}(2))$.
\end{proof}

\begin{proposition}[Steps $2$ through $n-3$: shifting window]\label{prop:shift}
For $2 \le j \le n - 3$,
\[
   \widetilde{W}_j = \mathrm{span}(\widetilde{u}_j),
   \qquad
   \widetilde{u}_j = A_j v_{(j,n-1)} + B_j v_{(j+1,n-1)}
                    + C_j v_{(j,n)} + D_j v_{(j+1,n)},
\]
with $A_j, B_j, C_j, D_j \in \mathbb{Q}(q)$ all nonzero.
\end{proposition}

\begin{proof}
By induction on $j$.

\paragraph{Base case $j = 2$.}  Apply $(T_2{+}1)$ to
$\widetilde{W}_1 = \mathrm{span}(w_{n-1}(2))$ where
$w_{n-1}(2) = \alpha v_{(2,n-1)} + \beta v_{(2,n)}$.

For $c \in \{n-1,n\}$ ($c \ge 4$): in $T_{(2,c)}$, $a = 2$ at $(1,2)$,
$b = 3$ at $(2,1)$ (smallest of $\{2,\ldots,n\}\setminus\{2,c\}$ is
$3$).  $T_2$ couples $v_{(2,c)}$ with $v_{(3,c)}$ (axial distance
$-2$), so $(T_2{+}1)v_{(2,c)} = \kappa(\alpha_2 v_{(2,c)} +
\beta_2 v_{(3,c)})$ with $\alpha_2,\beta_2,\kappa$ nonzero.

Hence $(T_2{+}1)w_{n-1}(2) = \kappa\bigl(\alpha\alpha_2 v_{(2,n-1)} +
\alpha\beta_2 v_{(3,n-1)} + \beta\alpha_2 v_{(2,n)} + \beta\beta_2 v_{(3,n)}\bigr)$,
all four coefficients nonzero.  This matches the form
$\widetilde{u}_2$ with $j = 2$, supported on
$\{(2,n-1),(3,n-1),(2,n),(3,n)\}$.

\paragraph{Inductive step $j-1 \to j$.}  Assume $\widetilde{W}_{j-1}$
is the span of $\widetilde{u}_{j-1}$ supported on
$\{(j-1,n-1),(j,n-1),(j-1,n),(j,n)\}$ with all four coefficients
nonzero.  Apply $(T_j{+}1)$ for $2 \le j \le n-3$.

\emph{(i) $v_{(j-1,c)}$ for $c \in \{n-1, n\}$.} $a = j-1, c \in \{n-1,n\}$.
Since $j \le n-3 < n-1 \le c$, $j$ and $j+1$ are both
$\le n-2 < n-1$, so neither equals $a = j-1$ nor $c$.  Hence $j, j+1$
are in column~$1$.

Position counting: column~$1$ entries are $\{1\}\cup\{2,\ldots,n\}\setminus\{j-1,c\}$,
sorted.  Elements less than $j$ in column~$1$ are
$\{1,\ldots,j-1\}\setminus\{j-1\} = \{1,\ldots,j-2\}$, totaling $j-1$
elements.  So $j$ is at row $j$, position $(j,1)$, content $1-j$.
Elements less than $j+1$ in column~$1$ are
$\{1,\ldots,j\}\setminus\{j-1\} = \{1,\ldots,j-2,j\}$, totaling $j-1$.
Wait: $\{1,\ldots,j\}$ has $j$ elements, removing one gives $j-1$ if
that one is in the set (and $j-1 \in \{1,\ldots,j\}$ iff $j \ge 2$).
So $j-1$ elements.  Hmm, this would put $j+1$ at row $j$, but $j$ is
already there.  Let me recount.

Number of column-$1$ elements strictly less than $j$: we need
elements of $\{1\}\cup\{2,\ldots,n\}\setminus\{j-1,c\}$ that are
strictly less than $j$.  Elements: $\{1,2,\ldots,j-1\}\setminus\{j-1\}
= \{1,\ldots,j-2\}$ (since $c > j$), total $j - 2$ elements (for
$j \ge 2$).  So $j$ is at row $j-1$, position $(j-1,1)$, content
$-(j-2) = 2-j$.

Elements strictly less than $j+1$ in column~$1$: $\{1,\ldots,j\}\setminus\{j-1\}
= \{1,\ldots,j-2,j\}$, total $j-1$ elements.  So $j+1$ is at row $j$,
position $(j,1)$, content $1-j$.

So $j$ at $(j-1,1)$, $j+1$ at $(j,1)$: same column adjacent.
$T_jv_{(j-1,c)} = -v_{(j-1,c)}$ for $c \in \{n-1,n\}$.
Hence $(T_j{+}1)v_{(j-1,c)} = 0$.

\emph{(ii) $v_{(j,c)}$ for $c \in \{n-1, n\}$.} $a = j$ at $(1,2)$,
content $1$.  $j+1$ in column~$1$ (not equal to $a = j$ or $c$, since
$j+1 \le n-2 < n-1 \le c$).

Position of $j+1$: number of column-$1$ entries less than $j+1$ is
$|\{1\}\cup\{2,\ldots,j\}\setminus\{j,c\}| = |\{1\}\cup\{2,\ldots,j-1\}|
= j-1$.  So $j+1$ at row $j$, position $(j,1)$, content $1-j$.

Different from $(1,2)$.  $T_j$-block partner: swap $j,j+1$ moves
$j+1$ to $(1,2)$, $j$ to $(j,1)$; new SYT $T_{(j+1,c)}$.  Axial
distance $1-j-1 = -j$.  $(T_j{+}1)v_{(j,c)} = \kappa_j(\alpha'v_{(j,c)}
+ \beta'v_{(j+1,c)})$ with $\kappa_j, \alpha', \beta'$ nonzero.

\emph{Combining.}
\begin{align*}
   (T_j{+}1)\,\widetilde{u}_{j-1}
   &= (T_j{+}1)\bigl(A_{j-1}v_{(j-1,n-1)} + B_{j-1}v_{(j,n-1)}
      + C_{j-1}v_{(j-1,n)} + D_{j-1}v_{(j,n)}\bigr) \\
   &= 0 + B_{j-1}\kappa_j(\alpha'v_{(j,n-1)} + \beta'v_{(j+1,n-1)})
      + 0 + D_{j-1}\kappa_j(\alpha'v_{(j,n)} + \beta'v_{(j+1,n)}) \\
   &= \kappa_j\bigl(B_{j-1}\alpha'\,v_{(j,n-1)} + B_{j-1}\beta'\,v_{(j+1,n-1)}
      + D_{j-1}\alpha'\,v_{(j,n)} + D_{j-1}\beta'\,v_{(j+1,n)}\bigr),
\end{align*}
which has all four coefficients nonzero ($B_{j-1}, D_{j-1}, \alpha',
\beta',\kappa_j$ all nonzero).  This is exactly $\widetilde{u}_j$
supported on $\{(j,n-1),(j+1,n-1),(j,n),(j+1,n)\}$.
\end{proof}

\begin{proposition}[Step $n-2$: terminal shape]\label{prop:terminal}
$\widetilde{W}_{n-2}$ is $1$-dimensional, supported on the $4$ SYTs
\[
   \{(n-3,n-2),\;(n-3,n-1),\;(n-2,n),\;(n-1,n)\},
\]
and contained in $E_1^-$.
\end{proposition}

\begin{proof}
Apply $(T_{n-2}{+}1)$ to $\widetilde{u}_{n-3}$ supported on
$\{(n-3,n-1),(n-2,n-1),(n-3,n),(n-2,n)\}$ (by Prop.~\ref{prop:shift}
with $j = n-3$, swapping the roles): write
$\widetilde{u}_{n-3} = A v_{(n-3,n-1)} + B v_{(n-2,n-1)}
+ C v_{(n-3,n)} + D v_{(n-2,n)}$ with $A,B,C,D \ne 0$.

\emph{(i) $v_{(n-3,n-1)}$.} $a = n-3$ at $(1,2)$, $c = n-1$ at $(2,2)$.
$n-2$: not in $\{a, c\}$.  Position: column~$1$ entries are
$\{1\}\cup\{2,\ldots,n\}\setminus\{n-3,n-1\} =
\{1,2,\ldots,n-4,n-2,n\}$.  Elements less than $n-2$ in this list:
$\{1,2,\ldots,n-4\}$, totaling $n-4$.  So $n-2$ at row $n-3$, position
$(n-3,1)$, content $-(n-4)$.  $n-1$ at $(2,2)$, content $0$.
Different.  $T_{n-2}$-block partner: swap $n-2, n-1$ moves $n-1$ to
$(n-3,1)$, $n-2$ to $(2,2)$; new SYT $T_{(n-3,n-2)}$.  Axial distance
$0 - (-(n-4)) = n-4$.
$(T_{n-2}{+}1)v_{(n-3,n-1)} = \kappa(\alpha'v_{(n-3,n-1)} +
\beta'v_{(n-3,n-2)})$ with $\kappa,\alpha',\beta' \ne 0$.

\emph{(ii) $v_{(n-2,n-1)}$.} $a = n-2, c = n-1$: same column~$2$
($a + 1 = c$).  $T_{n-2}v = -v$.  $(T_{n-2}{+}1)v = 0$.

\emph{(iii) $v_{(n-3,n)}$.} $a = n-3, c = n$.  $n-2,n-1$: both not in
$\{a,c\}$, so both in column~$1$.  Column~$1$ entries:
$\{1\}\cup\{2,\ldots,n\}\setminus\{n-3,n\} = \{1,2,\ldots,n-4,n-2,n-1\}$.
$n-2$ at row $n-3$, content $-(n-4)$; $n-1$ at row $n-2$, content
$-(n-3)$.  Same column adjacent.  $T_{n-2}v = -v$.
$(T_{n-2}{+}1)v = 0$.

\emph{(iv) $v_{(n-2,n)}$.} $a = n-2$ at $(1,2)$, content $1$.
$c = n$ at $(2,2)$.  $n-1$: column~$1$ (not equal to $a$ or $c$).
Column~$1$: $\{1\}\cup\{2,\ldots,n\}\setminus\{n-2,n\} =
\{1,2,\ldots,n-3,n-1\}$.  $n-1$ at row $n-2$, position $(n-2,1)$,
content $-(n-3)$.  Different from $(1,2)$.  $T_{n-2}$-block partner:
swap $n-2, n-1$ moves $n-1$ to $(1,2)$, $n-2$ to $(n-2,1)$; new SYT
$T_{(n-1,n)}$.  Axial distance $-(n-3)-1 = -(n-2)$.
$(T_{n-2}{+}1)v_{(n-2,n)} = \kappa'(\alpha''v_{(n-2,n)} +
\beta''v_{(n-1,n)})$ with $\kappa',\alpha'',\beta'' \ne 0$.

\emph{Combining.}
\begin{align*}
   (T_{n-2}{+}1)\,\widetilde{u}_{n-3}
   &= A\kappa(\alpha'v_{(n-3,n-1)} + \beta'v_{(n-3,n-2)}) + 0 + 0 \\
   &\quad + D\kappa'(\alpha''v_{(n-2,n)} + \beta''v_{(n-1,n)}).
\end{align*}
Four nonzero coefficients (since $A,D,\kappa,\kappa',\alpha',\alpha'',\beta',\beta''
\ne 0$); supports as claimed.

\emph{$E_1^-$ containment.}  Every support has minimum element
$\ge n-3 \ge 3$ (for $n \ge 6$).  By Lemma~\ref{lem:E1-22},
$2 \notin S$ implies $v_S \in E_1^-$.  Hence $\widetilde{u}_{n-2} \in E_1^-$.
\end{proof}

\begin{proof}[Proof of Theorem~\ref{thm:main}]
Combine Proposition~\ref{prop:phaseA} (Phase A endpoint
$W_n = E_{n-1}^+$, dim $n-3$),
Proposition~\ref{prop:step1} (rank collapses to $1$),
Proposition~\ref{prop:shift} (shifting window),
Proposition~\ref{prop:terminal} (terminal step).
$B_{n-1}V_\lambda = R'_{n-1}R'_n V_\lambda = \widetilde{W}_{n-2}$
is one-dimensional, supported on
$\{(n-3,n-2),(n-3,n-1),(n-2,n),(n-1,n)\}$, and contained in $E_1^-$.
\end{proof}

\section{Computational verification}

\begin{theorem}[Verification at $n=6,\ldots,9$]\label{thm:verify}
Theorem~\ref{thm:main} and the intermediate steps
(Propositions~\ref{prop:phaseA}, \ref{prop:step1}, \ref{prop:shift},
\ref{prop:terminal}) hold computationally at $n = 6, 7, 8, 9$.
\end{theorem}

\begin{proof}[Verification]
SymPy at $q = 7 \in \mathbb{Q}$ (generic for the seminormal basis).
At each step, the rank, support pattern, and $E_1^-$ containment match
the predicted values.  At $n=9$, $\widetilde{W}_{n-2}$ is supported
on $\{(6,7),(6,8),(7,9),(8,9)\}$, all with min element $6 \ge 3$,
hence in $E_1^-$.  Script:
\texttt{\textasciitilde/projects/scratch/2026-05-13-non-hook-22-1n/verify\_proof.py}.
The script verifies:
\begin{itemize}[leftmargin=2em,topsep=0.3em,itemsep=0.2em]
\item Phase A: at every $j \in \{1,\ldots,n-1\}$, rank of $W_j^{(n)}$
is exactly $n-3$, and the support matches the explicit prediction
from $\mathcal{B}_j$.
\item Phase B: at every $j \in \{1,\ldots,n-2\}$, rank of
$\widetilde{W}_j$ is exactly $1$, and the support matches the
prediction.  $E_1^-$ containment confirmed for $j \ge 3$ (and
specifically at $j = n-2$).
\end{itemize}
\end{proof}

\section{Discussion}\label{sec:discuss}

\subsection{Comparison with the hook results}

The structural pattern --- Phase A landing in $E_{n-1}^+$, Phase B
collapsing to $1$-dimensional via $(T_1{+}1)$, then shifting through
a window, then a terminal step --- closely parallels the May-12 proof
for the hook $\lambda = (3, 1^{n-3})$ (writeup
\texttt{2026-05-12-hook-j-formula.tex}).

The key difference is the terminal support pattern:
\begin{itemize}[leftmargin=2em,topsep=0.3em,itemsep=0.2em]
\item Hook $(3,1^{n-3})$: $5$ SYTs in support, namely
$\{(n-3,n-2),(n-3,n-1),(n-2,n-1),(n-2,n),(n-1,n)\}$.
\item Non-hook $(2,2,1^{n-4})$: $4$ SYTs in support, namely
$\{(n-3,n-2),(n-3,n-1),(n-2,n),(n-1,n)\}$.
\end{itemize}
The non-hook case has \emph{one fewer} SYT in the support: namely
$(n-2,n-1)$ is excluded.  This is because in $\lambda=(2,2,1^{n-4})$
the SYT $T_{(n-2,n-1)}$ has $(a,c) = (n-2,n-1)$ in column~$2$
together --- a same-column-$2$ configuration for $T_{n-2}$, which
$(T_{n-2}{+}1)$ kills.  In the hook case, the analogous SYT has
$\{n-2,n-1\}$ both in row~$1$, which $(T_{n-2}{+}1)$ acts on as
$(1+q)$ (same-row), hence the SYT $(n-2,n-1)$ \emph{survives} and
contributes to the support.

\subsection{Cardinality of supports}

We have now seen the following support sizes for $B_{\ell+1} V_\lambda$
in the rank-zero regime:
\begin{itemize}[leftmargin=2em,topsep=0.3em,itemsep=0.2em]
\item $\lambda = (2, 1^{n-2})$ ($k = 2$ hook): support size $2$.
\item $\lambda = (3, 1^{n-3})$ ($k = 3$ hook): support size $5 = C_3$.
\item $\lambda = (2, 2, 1^{n-4})$ (current paper): support size $4$.
\end{itemize}
The hook support sizes match Catalan numbers $C_k$
(\texttt{2026-05-12-hook-j-formula.tex} Conj.~3.1).  The non-hook
support size $4$ does not directly match any obvious Catalan-type
count.  This raises the question:

\begin{conjecture}[Support size for general non-hooks]\label{conj:nonhook}
For $\lambda \vdash n$ in the rank-zero regime, the support size of
the unique (up to scalar) spanning vector of $B_{\ell+1}V_\lambda$ is
governed by a $\kappa(\lambda)$-tuple of Catalan-like numbers, one per
length-preserving corner of $\lambda$.  For $\lambda = (2,2,1^{n-4})$
(one corner at row~$2$): the relevant ``Catalan-like'' number is
$\binom{2}{1} \cdot C_2 = 2 \cdot 2 = 4$, matching our count.
\end{conjecture}

The factor $\binom{2}{1}$ would account for the choice of column~$1$
vs.\ column~$2$ at the corner: a non-hook corner allows ``two ways
out'' compared to a hook.  This is purely speculative; verifying the
formula on more non-hook cases is the natural next step.

\subsection{Coverage of the rank-zero theorem}

Together with the May-11 and May-12 proofs, the rank-zero theorem
$\Pi^{S_n}|_{V_\lambda} = 0$ is now \emph{structurally proven} for
the following infinite families:
\begin{itemize}[leftmargin=2em,topsep=0.3em,itemsep=0.2em]
\item $\lambda = (1^n)$ (sign rep): trivial, since $\Pi^{S_n}$
contains $(T_i{+}1)$ for every $i$, and $(T_i{+}1)$ acts as $0$ on
the sign rep.
\item $\lambda = (2, 1^{n-2})$ (May 11): hook $k=2$ family.
\item $\lambda = (3, 1^{n-3})$ for $n \ge 6$ (May 12): hook $k=3$
family.
\item $\lambda = (2, 2, 1^{n-4})$ for $n \ge 6$ (current paper):
first non-hook family.
\end{itemize}
The remaining cases in the rank-zero regime (e.g., $(4, 1^{n-4})$
hook, $(3, 2, 1^{n-5})$, $(2, 2, 2, 1^{n-6})$, etc.) remain open
structurally.

\section{Honest assessment}\label{sec:honest}

\paragraph{What this paper achieves.}
\begin{enumerate}[leftmargin=2em,topsep=0.3em,itemsep=0.2em]
\item Theorem~\ref{thm:main} is proven structurally for every
$n \ge 6$.  The proof is a step-by-step computation in the seminormal
basis, with $\kappa$ and similar nonzero scalars tracked but not
computed explicitly.
\item Corollary~\ref{cor:rank-zero}: $\Pi^{S_n}|_{V_{(2,2,1^{n-4})}}
= 0$ for all $n \ge 6$, the first non-hook structural rank-zero
result.
\item The terminal support has size $4$, not $C_3 = 5$ as in the
hook $(3,1^{n-3})$ case.  The structural reason --- one SYT
$(n-2,n-1)$ is excluded due to the same-column-$2$ killing --- is
identified explicitly.
\item Conjecture~\ref{conj:nonhook}: a tentative Catalan-product
formula for support sizes of non-hooks, with $\binom{k}{j}\cdot C_j$
as the relevant counting unit per length-preserving corner.
\end{enumerate}

\paragraph{What this paper does \emph{not} achieve.}
\begin{enumerate}[leftmargin=2em,topsep=0.3em,itemsep=0.2em]
\item Conjecture~\ref{conj:nonhook} is verified at $n \le 9$ for
$\lambda = (2,2,1^{n-4})$ but unverified for other non-hooks.
\item The relationship between the support of $B_{\ell+1} V_\lambda$
and the SYTs of $\lambda^{>1} = (1, 1)$ (the partition obtained by
removing column~$1$) is suggestive but not made precise.
\item Phase A of the proof requires a delicate inductive argument
(Lemma~\ref{lem:transition}) that handles three sub-cases.  A more
``categorical'' or ``functorial'' version, possibly via branching
$V_\lambda \downarrow_{S_{n-1}}$ or via an interpolation/limit
argument, would be aesthetically more satisfying.
\end{enumerate}

\paragraph{Status.}  The structural rank-zero theorem now covers two
hook families and one non-hook family.  The path forward is either
(a) extending to more non-hook families
($\lambda = (2,2,2,1^{n-6})$ being the most natural test case), or
(b) finding a uniform proof covering all $\lambda$ in the rank-zero
regime.  The latter likely requires a deeper invariant
(representation-theoretic or geometric) that we have not yet
identified.

\end{document}
